\SourceLeaf{350}%
geschrieben, oder die Clisen von $\beta\gamma$, $\alpha\gamma$, $\alpha\beta$ mit $2\mathfrak{A}$, $2\mathfrak{B}$, $2\mathfrak{C}$ bezeichnet werden:
\[
\frac{4\cos\mathfrak{A}\cos\mathfrak{B}\cos\mathfrak{C}}{\dfrac{1}{n}+\dfrac{1}{m}} - \frac{4\sin\mathfrak{A}\sin\mathfrak{B}\sin\mathfrak{C}}{\dfrac{1}{n}-\dfrac{1}{m}}\, i.
\]
Es sind hierbei aber die Halbkreise, in welchen $\mathfrak{A}$, $\mathfrak{B}$, $\mathfrak{C}$ genommen werden, nur für zwei willkürlich, z.\,B. für $\mathfrak{A}$, $\mathfrak{B}$; für den dritten muss werden
\[
\gamma = \cos(\mathfrak{A}+\mathfrak{B}-\mathfrak{C}) + i\sin(\mathfrak{A}+\mathfrak{B}-\mathfrak{C}).
\]
\Needspace{12\baselineskip}Mit andern Worten, hätte man gemacht
\[
a = \cos a^{*} + i\sin a^{*}
\]
\[
\beta = \cos \beta^{*} + i\sin \beta^{*}
\]
\[
\gamma = \cos \gamma^{*} + i\sin \gamma^{*},
\]
so würde man zu setzen haben
\[
2\mathfrak{A} = \beta^{*}+\gamma^{*}
\]
\[
2\mathfrak{B} = a^{*}+\gamma^{*}
\]
\[
2\mathfrak{C} = a^{*}+\beta^{*}.
\]
\Needspace{6\baselineskip}Der Halbmesser des Kreises wird
\begingroup\small
\[
= \frac{\BareRoot{((mm+nn-2mn\cos 2\mathfrak{A})(mm+nn-2mn\cos 2\mathfrak{B})(mm+nn-2mn\cos 2\mathfrak{C}))}}{mm-nn}.
\]
\endgroup

\Needspace{11\baselineskip}\begin{center}[3.]\end{center}

Bei der \textsc{Newton}schen Construction hat die Differenz zweier Werthe von $t$ die Clise wie
\[
\frac{\dfrac{\beta\gamma}{\beta-\gamma}(b-c)\times\dfrac{\beta-\delta}{b-d}\times\dfrac{\gamma-\delta}{c-d}}{\beta\gamma\delta}
\left\{\frac{1}{\alpha\beta\gamma\delta-1}-\frac{1}{\alpha'\beta\gamma\delta-1}\right\}
\]
oder, weil $\beta-\delta$ die Clise [wie] $i\sqrt{\beta\delta}$ u.\,s.\,w.,
\[
\frac{(\beta-\delta)(\gamma-\delta)}{\delta(\beta-\gamma)}
\left\{\frac{(\alpha-\alpha')\beta\gamma\delta}{(\alpha\beta\gamma\delta-1).(\alpha'\beta\gamma\delta-1)}\right\}
\frac{b-c}{(b-d)(c-d)}
\]
oder, weil Clise von $\alpha-\alpha'$ wie $i\sqrt{\alpha\alpha'}$ u.\,s.\,w., wie
\[
\frac{b-c}{(b-d)(c-d)}.
\]
